\documentclass[12pt]{article}
\usepackage[utf8]{inputenc}
\usepackage[T1]{fontenc}
\usepackage[margin=1in]{geometry}
\usepackage{setspace}
\usepackage{xcolor}
\usepackage[colorlinks=true, urlcolor=blue!60!black, linkcolor=blue!60!black]{hyperref}
\setlength{\parindent}{2em}
\setlength{\parskip}{1.25em}
\renewcommand{\baselinestretch}{1.0}
\begin{document}
\begin{center}
{\LARGE \textit{Medell\'{i}n v. Texas}, 552 U.S. 491 (2008)\par}
\vspace{0.8em}
{\large Case Brief --- Four-Page Edition\par}
\end{center}
\vspace{1em}

\section{Memory Jogger}

The Court held 6--3 that neither the ICJ's \textit{Avena} judgment nor President Bush's memorandum directing state courts to give it effect was binding domestic law: a non-self-executing treaty obligation becomes enforceable in U.S.\ courts only through congressional implementing legislation (\href{https://supreme.justia.com/cases/federal/us/552/491/}{552 U.S.\ 491}).

\section{Detailed Case Facts}

In 1969 the United States ratified the Vienna Convention on Consular Relations, whose Article 36 requires authorities to tell arrested foreign nationals of their right to contact their consulate, and its Optional Protocol submitting Convention disputes to the International Court of Justice (ICJ). Under Article 94(1) of the U.N.\ Charter, each member ``undertakes to comply with the decision of the [ICJ] in any case to which it is a party.''

Jos\'{e} Ernesto Medell\'{i}n, a Mexican national who had lived in the United States since preschool, was arrested in 1993 for the gang rape and murder of two Houston teenagers. He waived his \textit{Miranda} rights and confessed, but officers never told him of his right to notify the Mexican consulate.

\section{Procedural History}

Medell\'{i}n was convicted and sentenced to death. He first raised the Convention claim in state postconviction proceedings, and the Texas courts, a federal district court, and the Fifth Circuit found it procedurally defaulted. The ICJ then held in \textit{Avena} (2004) that the United States had violated Article 36 as to 51 Mexican nationals, including Medell\'{i}n, and must provide ``review and reconsideration'' regardless of procedural default. After the Court granted certiorari in \textit{Medell\'{i}n I}, President Bush directed that the United States would comply ``by having State courts give effect to the decision in accordance with general principles of comity.'' The Court dismissed \textit{Medell\'{i}n I}, held in \textit{Sanchez-Llamas v.\ Oregon} (2006) that the Convention did not displace state default rules, and, after the Texas Court of Criminal Appeals dismissed a second habeas application as an abuse of the writ, again granted certiorari. It heard argument on October 10, 2007, and affirmed on March 25, 2008.

\section{Judicial Votes}

Affirmed, \textbf{6--3}.
\begin{itemize}
\item \textbf{Majority:} Chief Justice Roberts, joined by Justices Scalia, Kennedy, Thomas, and Alito.
\item \textbf{Concurring in the judgment:} Justice Stevens.
\item \textbf{Dissenting:} Justice Breyer, joined by Justices Souter and Ginsburg.
\end{itemize}

\section{Holding}

\textit{Avena} creates an international obligation, but neither the judgment nor Article 94 of the U.N.\ Charter is self-executing federal law enforceable in U.S.\ courts without implementing legislation. Under Category Three of Justice Jackson's \textit{Youngstown} framework, the President lacks unilateral authority to convert a non-self-executing treaty into domestic law that preempts state procedural default rules.

\section{Analysis of Opinions}

\textbf{Majority (Roberts).} \emph{Self-execution:} under \textit{Foster v.\ Neilson} (1829), a treaty is domestic law only if Congress implements it or its text shows it operates of its own force. The Optional Protocol merely grants ICJ jurisdiction; ``undertakes to comply'' in Article 94(1) commits the political branches to future action; Article 94(2) supplies a diplomatic remedy through the Security Council, where the United States has a veto; and only nations may be ICJ parties. The Executive's interpretation and the practice of all 47 Protocol signatories confirmed the point. \emph{Presidential power:} though the President's interests were ``plainly compelling,'' his authority ``must stem either from an act of Congress or from the Constitution itself.'' Making a non-self-executing treaty domestic law conflicted with the ratifying Senate's understanding, placing him at his ``lowest ebb''; that task belongs to Congress. \textit{Dames \& Moore} and \textit{Garamendi} rested on a long, congressionally acquiesced practice of settling civil claims, not ordering state courts to reopen criminal judgments, and the Take Care Clause did not apply because \textit{Avena} was not domestic ``Law.''

\textbf{Stevens (concurring in the judgment).} Stevens rejected any presumption against self-execution and thought the Convention self-executing, but read ``undertakes to comply'' as contemplating political action. Citing Oklahoma's voluntary compliance, he urged Texas to protect the Nation's honor by granting review.

\textbf{Breyer (dissenting).} Multilateral drafters rarely address domestic self-execution. Drawing on precedent since \textit{Ware v.\ Hylton}, Breyer proposed a practical, context-specific inquiry: the underlying Convention right is self-executing, the Protocol provides for ``compulsory settlement,'' reviewing prejudice is judicial work, and the President supported enforcement. He warned that some 70 treaties with ICJ clauses were jeopardized. After the decision, Congress passed no implementing law, and Texas executed Medell\'{i}n in August 2008.

\section{Examples: Future Applications}

\begin{enumerate}
\item \textbf{Same-Side (Not Domestic Law): Environmental Tribunal Directive.} A treaty in which states ``undertake to give effect'' to tribunal rulings, with diplomatic enforcement only, is non-self-executing, and a presidential directive ordering California courts to enjoin a zoning law falls in Category Three.
\item \textbf{Same-Side (Not Domestic Law): Non-Self-Executing Declaration.} A human-rights treaty ratified subject to a Senate declaration of non-self-execution gives a state prisoner no rule of decision.
\item \textbf{Opposite-Side (Binding): Implementing Legislation or Explicit Text.} A statute like the ICSID enforcement law, 22 U.S.C.\ \S\ 1650a, or treaty text directing domestic courts to enforce tribunal orders, makes the obligation supreme federal law.
\item \textbf{Opposite-Side (Valid): Claims-Settlement Agreement.} An executive agreement settling U.S.\ citizens' civil claims against a foreign sovereign falls within the tradition upheld in \textit{Dames \& Moore} and \textit{Garamendi}.
\item \textbf{Fence-Sitter: Section 1983 Damages for Missing Consular Notice.} A suit resting on Article 36 itself avoids the \textit{Avena} problem, but \textit{Sanchez-Llamas} only assumed individual rights exist, and the Seventh and Ninth Circuits have split on damages.
\end{enumerate}

\section{Critique}

\textbf{Federalism versus one voice.} The Supremacy Clause was adopted to keep individual states from putting the Nation in breach of its treaties, yet \textit{Medell\'{i}n}'s holdings let a single state's procedural rule cause an acknowledged international-law violation unless Congress acts.

\textbf{Textual clarity in multilateral treaties.} As Breyer emphasized, foreign negotiators rarely draft clauses addressing U.S.\ self-execution, so critics argue that \textit{Medell\'{i}n} effectively creates a presumption against self-execution despite disclaiming any need for ``talismanic words.''

\section{Key Quotations}

\begin{quote}
``In sum, while treaties `may comprise international commitments . . . they are not domestic law unless Congress has either enacted implementing statutes or the treaty itself conveys an intention that it be ``self-executing'' and is ratified on these terms.' ''
\end{quote}
\noindent --- Chief Justice Roberts, 552 U.S.\ at 505.

\begin{quote}
``The President has an array of political and diplomatic means available to enforce international obligations, but unilaterally converting a non-self-executing treaty into a self-executing one is not among them. The responsibility for transforming an international obligation arising from a non-self-executing treaty into domestic law falls to Congress.''
\end{quote}
\noindent --- Chief Justice Roberts, 552 U.S.\ at 525--26.

\begin{quote}
``When the President asserts the power to `enforce' a non-self-executing treaty by unilaterally creating domestic law, he acts in conflict with the implicit understanding of the ratifying Senate. His assertion of authority, insofar as it is based on the pertinent non-self-executing treaties, is therefore within Justice Jackson's third category, not the first or even the second.''
\end{quote}
\noindent --- Chief Justice Roberts, 552 U.S.\ at 527.

\end{document}
